%% ==================== AI 工具使用声明 模板 (v1.5.6, BZD 数模社 2026 规范) ====================
%%
%% 借鉴 BZD 数模社 论文末尾模板: 二者择一 (未使用 AI / 使用 AI)
%% 必须在参考文献之前, AI 工具不列入参考文献, 详细情况见支撑材料《AI工具使用详情.pdf》
%%
%% 跑题用户必须:
%% 1. 删掉所有 % 注释
%% 2. 二者择一: 删除其中一段, 保留另一段
%% 3. 替换【】占位符为自己的实际内容
%% 4. 删 下面的 \AISlot{...} 行 (check 15 抓这个 inline 【】)

\subsection*{AI 工具使用声明}

% ====== 二者择一: 跑题用户必须删掉其中一段, 保留另一段 ======

% ===== 选项 (1) 未使用 AI 工具: 保留下面这行, 删掉 "选项 (2)" =====
% 本参赛队在竞赛过程中未使用任何 AI 工具。

% ===== 选项 (2) 使用 AI 工具: 删掉 "选项 (1)" 上面这行, 保留下面这段, 并填实际用途 =====
%% ⚠️ §9 论文内 = 1 句话简短声明, 只填 3-5 项 **简要** 用途 (e.g. 读题 / 代码 / 公式 / 文档 / 可视化)
%%    **不要在这里展开 4 节详写**, 详细 5+ 项看 支撑材料/AI工具使用详情.pdf (4 节)
本参赛队在竞赛过程中使用了 AI 工具, 主要用于【简要填 3-5 项实际用途, e.g. 读题 / 代码 (LP + CBC) / 公式推导 / 文档润色 / 数据可视化】, 详细使用情况见支撑材料《AI 工具使用详情.pdf》。

% 跑题后必须删下面这一行 (check 15 占位符自检)
\AISlot{【这里填实际使用情况, 跑题后删这一行】}

% AI 工具不能列入参考文献! ChatGPT / DeepSeek / Claude / Copilot 等条目禁止.
% AI 推荐的文献只列原始论文 (经人工核验), 不引用 AI.
% 声明应与支撑材料《AI工具使用详情.pdf》内容一致.
